% !TeX program = lualatex
% Standalone notebook; compile twice: lualatex 32.tex
% Research batch 39 down to 25, 27 September 2026.
% Original section preserved verbatim except for marked insertion blocks.
\documentclass[11pt,a4paper]{article}
\usepackage[margin=24mm,headheight=14pt]{geometry}
\usepackage{fontspec}
\setmainfont{Latin Modern Roman}
\setsansfont{Latin Modern Sans}
\setmonofont{Latin Modern Mono}
\usepackage{amsmath,amssymb,mathtools}
\usepackage{unicode-math}
\setmathfont{Latin Modern Math}
\usepackage{microtype}
\usepackage{enumitem,xurl,needspace}
\usepackage[dvipsnames]{xcolor}
\usepackage{hyperref}
\hypersetup{unicode=true,colorlinks=true,linkcolor=MidnightBlue,urlcolor=MidnightBlue,
 pdftitle={Research notebook: Problem 32},
 pdfauthor={Open Problems Math research notebook},bookmarksnumbered=true}
\usepackage{bookmark,titlesec,fancyhdr}
\titleformat{\section}{\Large\bfseries\raggedright}{\thesection}{0.7em}{}
\pagestyle{fancy}\fancyhf{}
\fancyhead[L]{\small\sffamily Open Problems Math\enspace|\enspace Research notebook}
\fancyhead[R]{\small\sffamily Problem 32}
\fancyfoot[L]{\small\sffamily 27 September 2026}
\fancyfoot[R]{\small\sffamily \thepage}
\setlength{\parindent}{0pt}
\setlength{\parskip}{5pt plus 1pt}
\setlength{\emergencystretch}{3em}
\setcounter{secnumdepth}{1}
\setcounter{tocdepth}{1}
\setlist[itemize]{leftmargin=3em,itemsep=5pt,topsep=4pt,parsep=0pt}
\allowdisplaybreaks
\newcommand{\N}{\mathbb{N}}
\newcommand{\Z}{\mathbb{Z}}
\newcommand{\Q}{\mathbb{Q}}
\newcommand{\R}{\mathbb{R}}
\newcommand{\C}{\mathbb{C}}
\newcommand{\F}{\mathbb{F}}
\newcommand{\A}{\mathbb{A}}
\newcommand{\PP}{\mathbb P}
\newcommand{\E}{\mathbb{E}}
\newcommand{\eps}{\varepsilon}
\newcommand{\dd}{\,\mathrm d}
\newcommand{\sref}[2]{\hyperlink{source:#1:#2}{[S#2]}}
\newcommand{\eref}[2]{\hyperlink{extra:#1:#2}{[E#2]}}
\newif\ifshowcatalogue
\showcataloguetrue
\newcommand{\cataloguescope}[1]{\ifshowcatalogue
 \paragraph{Exact scope in the supplied catalogue.}
 \begin{quote}\small #1\end{quote}\fi}
\numberwithin{equation}{section}
\begin{document}
\setcounter{section}{31}
% BEGIN PRESERVED SECTION
% ================================================================
% problemId: problem.valiants-algebraic-p-versus-np-conjecture-vp-vnp
% source releaseRank: 32; source releaseStatus: open
% exactTarget SHA-256: 4725cf7659e3460efd4c63f6e286ffa30d833740fea013ba0a69afa251ed2be2
\clearpage
\section{\texorpdfstring{Valiants Algebraic P versus NP Conjecture (VP != VNP)}{Valiants Algebraic P versus NP Conjecture (VP != VNP)}}\label{32}
\subsection{Definitions and mathematical statement}
For variables $x_{ij}$ over a characteristic-zero field $K$, define
\[
 \operatorname{per}_n(x)=\sum_{\sigma\in S_n}\prod_{i=1}^n x_{i,\sigma(i)}.
\]
An algebraic circuit computes a polynomial using field constants, variables, addition, and multiplication. The class $\mathsf{VP}$ has polynomially bounded numbers of variables, degrees, and circuit sizes; $\mathsf{VNP}$ consists of polynomially long Boolean sums of such efficiently computable polynomials. The selected claim is that no polynomial bounds the arithmetic circuit sizes of $\operatorname{per}_n$, equivalently $\mathsf{VP}\ne\mathsf{VNP}$ in characteristic zero. This is a nonuniform algebraic-operation model, not a bit-complexity assertion.
\par\smallskip\textit{Formulation and concept references:} \sref{32}{1}.
\cataloguescope{Prove that the permanent polynomial family cannot be computed by polynomial-size algebraic circuits over fields of characteristic zero, equivalently that VP is not equal to VNP.}
\subsection{Short English statement}
Can one prove that the permanent requires more than polynomially many arithmetic operations in general algebraic circuits?
% BEGIN RESEARCH 32
\subsection{Research attempt}

\paragraph{Current frontier.}
The supplied 2026 formulation paper proves separations for symmetry-constrained
algebraic models and completeness results for homomorphism polynomials. It
explicitly distinguishes these from unrestricted $\mathsf{VP}$ versus
$\mathsf{VNP}$. Output symmetry is not an assumption on every circuit
computing that output. \sref{32}{1}, \sref{32}{2}.
Narayanan's July 2026 manuscript uses sumset expansion to construct elusive
functions and improve certain low-depth bounds. Its abstract still identifies
the stronger explicit constructions needed for Valiant's conjecture as
missing. \sref{32}{3}, \eref{32}{1}. The relevant definitions and statements
were inspected; no full proof audit of the recent manuscripts is claimed.
Partial-derivative methods have a long-established role in restricted
arithmetic lower bounds. \eref{32}{2}.

\paragraph{Attack route.}
A direct route would attach a large linear-algebraic invariant to the
permanent and prove that it grows only polynomially in unrestricted circuit
size. A second route would use the permanent's symmetry to replace a
hypothetical small circuit by a small symmetric one, then invoke the known
symmetric separation. I test the first route with exact derivative and
row-cut profiles and use the second as an adversarial comparison. The goal
is to identify an invariant that sees cancellation and circuit sharing,
rather than merely count monomials or minors.

\paragraph{Attempt.}
For a homogeneous polynomial $f$ over a characteristic-zero field $K$, write
\[
 D_k(f)=\dim_K\operatorname{span}\{\partial^\alpha f:|\alpha|=k\}.
\]
The following calculation gives the complete profile, not just one
asymptotic lower bound.

\textbf{Lemma 32.1 (identical derivative dimensions).}
For $0\le k\le n$,
\begin{equation}\label{r32:derivatives}
 D_k(\operatorname{per}_n)=D_k(\det_n)=\binom nk^2.
\end{equation}
For $k>n$ both quantities are zero.

\emph{Proof.} Differentiating a permanent monomial at $k$ entries gives
zero unless the entries lie in distinct rows and distinct columns. In the
nonzero case, the result is the permanent of the complementary
$(n-k)\times(n-k)$ submatrix. Every pair of deleted row and column sets
occurs, giving at most $\binom nk^2$ polynomials. They are linearly independent:
the variables appearing in any monomial record exactly the remaining row
and column sets, so polynomials attached to different pairs have disjoint
monomial supports. Each is nonzero. The same argument for the determinant
gives signed complementary determinants, still nonzero with disjoint
supports. At $k=n$ all nonzero derivatives are constants and the formula
gives one.\hfill$\square$

The full profile is therefore incapable of distinguishing the two
polynomials. This is a sharper stress test than simply finding a weak bound
at one derivative order. In particular, no argument using only this profile
and a bound valid for \emph{all} small circuits can certify that the permanent
requires larger circuits than the known polynomial-size determinant circuit.
For clarity, the latter fact can be exhibited directly in this model.

Let $p_i=\operatorname{tr}(X^i)$, $e_0=1$, and recursively compute
\begin{equation}\label{r32:newton}
 e_j=\frac1j\sum_{i=1}^j(-1)^{i-1}e_{j-i}p_i
 \quad(1\le j\le n).
\end{equation}
The formal identity
$\det(I+tX)'/\det(I+tX)=\sum_{i\ge1}(-1)^{i-1}p_i t^{i-1}$
proves that $e_n=\det X$: it follows either by formal differentiation of the
determinant or by triangularizing over an algebraic closure and comparing
symmetric polynomials. Sequential matrix multiplication computes all $X^i$
with $O(n^4)$ arithmetic operations; the recurrence adds $O(n^2)$.
The divisions by $j$ are multiplications by allowed constants $1/j\in K$,
not division gates or variable-dependent denominators. Thus this is a
valid polynomial-size arithmetic circuit in every characteristic-zero field.
Its derivative dimensions near $k=n/2$ are nevertheless exponential in $n$.

The basic product inequalities explain the mismatch:
\[
 D_k(f+g)\le D_k(f)+D_k(g),\qquad
 D_k(fg)\le\sum_{i=0}^kD_i(f)D_{k-i}(g).
\]
Leibniz's rule proves the second. Even the $O(n)$-size circuit
$f=x_1\cdots x_n$ has $D_k(f)=\binom nk$. Exponential derivative-space
growth is compatible with a linear number of multiplication gates. The
missing invariant cannot be supplied by treating that growth as though it
were additive at product gates.

\medskip
\textbf{A restricted lower bound that survives the calculation.}
Split the variables into $Y$, the first $k$ rows, and $Z$, the other rows,
and let $\operatorname{rank}_{Y|Z}(f)$ be the rank of the matrix of coefficients
of monomials $u(Y)v(Z)$. All such matrices are finite for a fixed polynomial.

\textbf{Lemma 32.2 (exact row-cut rank).}
For $0\le k\le n$,
\begin{equation}\label{r32:cut}
 \operatorname{rank}_{Y|Z}(\operatorname{per}_n)
 =\operatorname{rank}_{Y|Z}(\det_n)=\binom nk.
\end{equation}
Consequently a decomposition
$f(Y,Z)=\sum_{j=1}^r a_j(Y)b_j(Z)$ of either polynomial requires
$r\ge\binom nk$, and that value is attainable.

\emph{Proof.} Group permutation monomials by the set $S$ of columns used
by the first $k$ rows. For the permanent this yields
\[
 \operatorname{per}X=\sum_{|S|=k}
     \operatorname{per}X_{[k],S}\,
     \operatorname{per}X_{[n]\setminus[k],S^c}.
\]
Laplace expansion gives the determinant version with a sign depending on
$S$. Distinct $S$ have disjoint supports on \emph{both} sides of the
coefficient matrix. Thus the matrix is a disjoint union of exactly
$\binom nk$ nonzero rank-one blocks. Each separated product contributes
at most rank one, proving both bounds.\hfill$\square$

At a balanced cut this is exponential, but it is a bound for a specified
separated-product representation. The determinant circuit above is free
to mix both row groups at internal gates, so there is no contradiction.
To use this route for the target one would need a polynomial-overhead
conversion of every arithmetic circuit into a representation respecting
an appropriate cut. Lemma~32.2 and the determinant circuit refute such a
conversion for this fixed row-cut model. A more flexible decomposition
would require a different invariant and proof.

\medskip
\textbf{Testing symmetrization as an escape.}
The permanent and determinant are both invariant under simultaneously
permuting rows and columns. Given any circuit for an invariant polynomial,
one can certainly make copies indexed by the group and average their
outputs. In characteristic zero the normalization by the group order is
legal. But for the group $S_n$ this straightforward construction uses
$n!$ copies; it gives no polynomial overhead. The fact that the final
polynomial is invariant does not certify a small group action on its
internal gates. This is exactly where the model distinction in
\sref{32}{1} and \sref{32}{2} matters. The observation does not prove that
every clever symmetrization must be expensive; it identifies the missing
lemma in this particular attempted transfer of their result.

\paragraph{Stress test.}
The derivative independence argument uses monomial supports, not numerical
values at sample matrices, so it is valid over every characteristic-zero
field and does not confuse generic rank with evaluated rank. Nonzero signs
in determinant minors do not change their supports. Arbitrary cancellations
and reused subcircuits remain allowed in the target; they have not been
replaced by monotone, formula, symmetric, or fixed-depth computation.

The accompanying exact script constructs the determinant and permanent for
$n\le4$, forms their derivative coefficient vectors, and checks the ranks
over $\Q$. It checks the row-cut ranks and Newton recurrence on small
integer matrices as independent implementation tests. These finite cases
cannot prove a superpolynomial asymptotic bound. Nor would a Boolean
bit-complexity argument automatically apply to unrestricted field constants
in this nonuniform arithmetic model.

\Needspace{8\baselineskip}
\paragraph{Outcome.}
\textbf{Outcome: Exploratory approach with unresolved gap.}

The attempt establishes exact derivative and cut-rank profiles and a sharp
lower bound for a separated-product subclass. Comparing them with an
explicit polynomial-size determinant circuit defeats the proposed transfer
to general circuits. The remaining task is an invariant, or a controlled
normal-form theorem, that detects the permanent's difficulty without also
misclassifying the determinant. No such theorem is supplied; the lemmas
are elementary structural checks and are not claimed to be historically
new or to separate $\mathsf{VP}$ from $\mathsf{VNP}$.

% END RESEARCH 32

\subsection{Sources}
\begin{itemize}\raggedright
\item[\textnormal{[S1]}]\hypertarget{source:32:1}{} Dwivedi, Pago, and Seppelt, Lower Bounds in Algebraic Complexity via Symmetry and Homomorphism Polynomials, STOC 2026, Introduction.
\url{https://www.prateekdwivedi.in/papers/symACT-confversion.pdf}.
{\small\textit{Catalogue formulation source.}}
\item[\textnormal{[S2]}]\hypertarget{source:32:2}{} Symmetric Algebraic Circuits and Homomorphism Polynomials.
\url{https://drops.dagstuhl.de/storage/00lipics/lipics-vol362-itcs2026/LIPIcs.ITCS.2026.46/LIPIcs.ITCS.2026.46.pdf}.
{\small\textit{Catalogue research/status reference; not a proof endorsement.}}
\item[\textnormal{[S3]}]\hypertarget{source:32:3}{} Arithmetic circuit lower bounds from sumset expansion.
\url{https://eccc.weizmann.ac.il/report/2026/138/}.
{\small\textit{Catalogue research/status reference; not a proof endorsement.}}
% BEGIN ADDED REFERENCES 32
\item[\textnormal{[E1]}]\hypertarget{extra:32:1}{} Anand Kumar Narayanan,
\emph{Arithmetic circuit lower bounds from sumset expansion},
arXiv:2607.15848 (2026), abstract and scope of the elusive-function
constructions; also ECCC TR26-138.
\url{https://arxiv.org/abs/2607.15848}.
\url{https://eccc.weizmann.ac.il/report/2026/138/}.
{\small\textit{Recent manuscript supplement to [S3], checked
27 September 2026; no unrestricted separation inferred.}}
\item[\textnormal{[E2]}]\hypertarget{extra:32:2}{} Noam Nisan and
Avi Wigderson, \emph{Lower Bounds on Arithmetic Circuits via Partial
Derivatives}, BRICS RS-95-43, preliminary version (1995).
\url{https://www.brics.dk/RS/95/43/index.html}.
{\small\textit{Foundational method and restricted-model scope; checked
27 September 2026. The notebook calculations are proved separately.}}

% END ADDED REFERENCES 32
\end{itemize}

% END PRESERVED SECTION
\end{document}
